\documentclass[UTF8,a4paper,11pt,openany]{ctexbook}
\usepackage{../common/statlearnbook}
\SLSetBookMetadata{第 5 册}{统计学习方法讲义 v2.6.0}{采样方法、主题模型与图排序}
\begin{document}
\setcounter{page}{708}
\setcounter{chapter}{30}
\setcounter{figure}{3}
\pagestyle{headings}
\chaptermark{马尔可夫链基础}
\sectionmark{遍历定理与可逆性}
\textbf{几何解释。}所有$n$维概率分布组成概率单纯形，映射$\rho\mapsto\rho A$把该单纯形映回自身。平稳分布是这个仿射映射的固定点。图\ref{fig:V5-C01-stationary-fixed-point}在两状态概率区间上画出$r\mapsto0.2+0.5r$，固定点$r=0.4$对应$\rho=(0.4,0.6)$。
% v2.3.1 figure UID: FIG-P556-01
% Proposed post-v2.3.1 polish for FIG-P556-01: protect key-label size and stroke hierarchy.
\tikzset{slfig-FIG-P556-01/.style={every path/.append style={line cap=round,line join=round}}}
\pgfplotsset{
  slfig-FIG-P556-01-axis/.style={
    tick label style={font=\fontsize{8.7pt}{10.2pt}\selectfont},
    label style={font=\fontsize{9.4pt}{11.2pt}\selectfont},
    clip=false,
  },
  slfig-FIG-P556-01-reference/.style={draw=SLGray,line width=.58pt,densely dashed},
  slfig-FIG-P556-01-map/.style={draw=SLTeal,line width=1.02pt},
  slfig-FIG-P556-01-cobweb-low/.style={draw=SLBlue,line width=.72pt},
  slfig-FIG-P556-01-cobweb-high/.style={draw=SLGold,line width=.72pt,densely dashed},
}
\begin{figure}[htbp]
\centering
\begin{tikzpicture}[slfig-FIG-P556-01]
\begin{axis}[slfig-FIG-P556-01-axis,
  width=9.2cm,height=7.3cm,
  xmin=0,xmax=1,ymin=0,ymax=1,
  axis equal image,axis lines=left,
  xlabel={$r_t$},ylabel={$r_{t+1}$},
  xtick={0,.2,.4,.6,.8,1},ytick={0,.2,.4,.6,.8,1},
  clip=false,
]
  \addplot[slfig-FIG-P556-01-reference,domain=0:1,samples=201] {x}
    node[pos=.88,above left,font=\fontsize{9.3pt}{11.0pt}\selectfont] {$y=r$};
  \addplot[slfig-FIG-P556-01-map,domain=0:1,samples=201] {.2+.5*x}
    node[pos=.94,above,yshift=1.6mm,font=\fontsize{9.3pt}{11.0pt}\selectfont,text=SLTeal] {$y=0.2+0.5r$};
  \addplot[only marks,mark=*,mark size=2.7pt,draw=SLGold,fill=SLGold]
    coordinates {(.4,.4)} node[above left,xshift=-.8mm,yshift=.8mm,font=\fontsize{9.4pt}{11.2pt}\selectfont,text=SLGold] {$r^*=0.4$};

  % r_0=0.05：实线 cobweb。
  \addplot[slfig-FIG-P556-01-cobweb-low] coordinates {
    (.05,0)(.05,.225)(.225,.225)(.225,.3125)(.3125,.3125)
    (.3125,.35625)(.35625,.35625)(.35625,.378125)(.378125,.378125)};
  \addplot[only marks,mark=square*,mark size=2.1pt,draw=SLBlue,fill=SLBlue]
    coordinates {(.05,0)};
  \node[anchor=west,yshift=-.6mm,font=\fontsize{9.2pt}{11.0pt}\selectfont,text=SLBlue] at (axis cs:.08,.075) {$r_0=.05$};

  % r_0=0.90：虚线 cobweb。
  \addplot[slfig-FIG-P556-01-cobweb-high] coordinates {
    (.90,0)(.90,.65)(.65,.65)(.65,.525)(.525,.525)
    (.525,.4625)(.4625,.4625)(.4625,.43125)(.43125,.43125)};
  \addplot[only marks,mark=triangle*,mark size=2.5pt,draw=SLGold,fill=SLGold]
    coordinates {(.90,0)};
  \node[anchor=east,yshift=-1.2mm,font=\fontsize{9.2pt}{11.0pt}\selectfont,text=SLGold] at (axis cs:.87,.075) {$r_0=.90$};
  \node[anchor=east,draw=SLRuleGray,line width=.65pt,rounded corners=2pt,fill=white,inner sep=2.2pt,align=center,font=\fontsize{9.2pt}{11.0pt}\selectfont]
    at (axis cs:.74,.18) {斜率 $0.5<1$\\压缩映射};
\end{axis}
\end{tikzpicture}
\caption{两状态概率单纯形上的平稳固定点：映射$r\mapsto0.2+0.5r$把任意初值逐步拉向$r_\star=0.4$}\label{fig:V5-C01-stationary-fixed-point}
\end{figure}

\textbf{概率解释。}不可约性回答“能否在状态之间往返”，周期性回答“回来的时刻是否被某个公因数锁定”，“正常返”回答“平均要等多久才能回来”，平稳性回答“总体概率质量是否保持不变”。四者回答的问题不同，不能用其中一个术语替代另一个。
\end{document}

